\section{大域コートクエリモデル}
\label{sec:model}

レンダリングパイプラインは、標準的なキーポイントヒートマップ損失だけでは学習できないカメラ仮説を
供給する。そこで、\cref{fig:model}に示すように、モデルを大域幾何エンコーダと、
意図的に軽量化した密予測ヘッドに分離する。

\begin{figure}[t]
  \centering
  \resizebox{\linewidth}{!}{\begin{tikzpicture}[
  font=\sffamily\footnotesize,
  module/.style={draw, rounded corners=2pt, align=center, minimum height=9mm,
    inner xsep=4pt, inner ysep=3pt, line width=0.55pt},
  patch/.style={draw=gsblue, fill=gsblue!16, rounded corners=1pt,
    minimum width=3.2mm, minimum height=5.0mm, inner sep=0pt},
  choice/.style={draw=courtgreen!75!black, fill=white, rounded corners=1pt,
    minimum width=22mm, minimum height=4.2mm, inner sep=1pt,
    font=\sffamily\scriptsize},
  head/.style={draw=gsblue!75!black, fill=white, rounded corners=1pt,
    minimum width=20mm, minimum height=4.2mm, inner sep=1pt,
    font=\sffamily\scriptsize\bfseries},
  arrow/.style={-{Latex[length=2.1mm,width=1.35mm]}, semithick,
    draw=black!72},
  taparrow/.style={-{Latex[length=2.1mm,width=1.35mm]}, semithick,
    draw=gsblue!85!black},
  queryarrow/.style={-{Latex[length=2.1mm,width=1.35mm]}, semithick,
    draw=poseorange!90!black},
  label/.style={font=\sffamily\scriptsize, text=black!72, align=center}
]
  \node[font=\sffamily\footnotesize\bfseries, text=gsblue!85!black, anchor=west]
    at (0,5.72) {(1) Patch tokens};
  \node[font=\sffamily\footnotesize\bfseries, text=poseorange!90!black, anchor=west]
    at (7.95,5.72) {(2) Global reasoning};
  \node[font=\sffamily\footnotesize\bfseries, text=courtgreen!85!black, anchor=west]
    at (11.55,5.72) {(3) Pose + dense readouts};

  \node[module, draw=gsblue, fill=gsblue!8, minimum width=16mm,
    minimum height=14mm] (rgb) at (0.85,4.00) {};
  \fill[courtgreen!88, rounded corners=0.7pt]
    ($(rgb.center)+(-0.57,0.01)$) rectangle ($(rgb.center)+(0.57,0.47)$);
  \draw[white, line width=0.45pt]
    ($(rgb.center)+(-0.48,0.08)$) rectangle ($(rgb.center)+(0.48,0.40)$);
  \draw[white, line width=0.4pt]
    ($(rgb.center)+(0,0.08)$) -- ($(rgb.center)+(0,0.40)$)
    ($(rgb.center)+(-0.48,0.24)$) -- ($(rgb.center)+(0.48,0.24)$);
  \node[font=\sffamily\scriptsize\bfseries] at ($(rgb.center)+(0,-0.40)$)
    {RGB frame};

  \node[module, draw=gsblue, fill=gsblue!10, minimum width=18mm,
    minimum height=12mm] (dino) at (2.75,4.00)
    {{\bfseries DINOv3} ViT\\[-0.15em]\scriptsize patch tokens only};
  \node[module, draw=poseorange, fill=poseorange!9, minimum width=15mm,
    minimum height=11mm] (projection) at (4.67,4.00)
    {Patch\\projection\\[-0.2em]\scriptsize $C\!\rightarrow\!d$};

  \node[module, draw=gsblue, fill=gsblue!5, minimum width=21mm,
    minimum height=12mm] (sequence) at (6.75,4.00) {};
  \node[patch] at ($(sequence.center)+(-0.60,0.16)$) {};
  \node[patch] at ($(sequence.center)+(-0.20,0.16)$) {};
  \node[patch] at ($(sequence.center)+(0.20,0.16)$) {};
  \node[font=\sffamily\scriptsize, text=gsblue!80!black]
    at ($(sequence.center)+(0.52,0.16)$) {$\cdots$};
  \node[label] at ($(sequence.center)+(0,-0.30)$)
    {$N=H_pW_p$ projected patches};

  \node[module, draw=poseorange, fill=poseorange!10, minimum width=22mm,
    minimum height=9mm] (query) at (6.75,1.72)
    {one learned court query\\[-0.15em]\scriptsize positionless: no RoPE};

  \node[module, draw=poseorange!35, fill=poseorange!3, minimum width=24mm,
    minimum height=14mm] at (9.58,4.18) {};
  \node[module, draw=poseorange!55, fill=poseorange!5, minimum width=24mm,
    minimum height=14mm] at (9.48,4.09) {};
  \node[module, draw=poseorange, fill=poseorange!11, minimum width=24mm,
    minimum height=14mm] (task) at (9.38,4.00)
    {Global task Transformer\\[-0.10em]
     \scriptsize shared self-attention + MLP\\[-0.10em]
     \scriptsize 2D RoPE on patch $Q,K$ only};

  \draw[arrow] (rgb.east) -- (dino.west);
  \draw[arrow] (dino.east) -- (projection.west);
  \draw[arrow] (projection.east) -- (sequence.west);
  \draw[arrow] (sequence.east) -- (task.west);
  \draw[queryarrow] (query.east) -| ($(task.south west)+(0.28,0)$);

  \node[module, draw=poseorange, fill=poseorange!10, text width=18mm,
    minimum height=9mm] (posehead) at (12.10,4.52)
    {final query\\[-0.15em]\scriptsize MLP pose head};
  \node[module, draw=poseorange, fill=poseorange!6, text width=32mm,
    minimum height=12mm] (poseout) at (15.15,4.52)
    {\bfseries raw pose10D\\[-0.10em]
     \scriptsize camera centre $\hat{\mathbf C}$ (3)\\[-0.10em]
     \scriptsize row--6D $\hat{\mathbf a}$ (6) $\;\vert\;$
     $\widehat{\log f}$ (1)};

  \node[choice] (linear) at (12.18,2.18) {Linear: $1{\times}1$ + resize};
  \node[choice] (progressive) at (12.18,1.70) {Progressive: upsample + DW conv};
  \node[choice] (dpt) at (12.18,1.22) {DPT: reassemble + fuse taps};
  \begin{scope}[on background layer]
    \node[draw=courtgreen, fill=courtgreen!8, rounded corners=2pt,
      fit=(linear)(progressive)(dpt), inner xsep=2.5mm, inner ysep=2.2mm]
      (decoders) {};
  \end{scope}
  \node[font=\sffamily\scriptsize\bfseries, text=courtgreen!85!black]
    at ($(decoders.north)+(0,0.23)$) {lightweight decoder (choose one)};

  \node[head] (kphead) at (15.28,2.18) {KP14};
  \node[head] (linehead) at (15.28,1.70) {LINE};
  \node[head] (seghead) at (15.28,1.22) {SEG};
  \begin{scope}[on background layer]
    \node[draw=gsblue, fill=gsblue!7, rounded corners=2pt,
      fit=(kphead)(linehead)(seghead), inner xsep=2.5mm, inner ysep=2.2mm]
      (heads) {};
  \end{scope}
  \node[font=\sffamily\scriptsize\bfseries, text=gsblue!85!black]
    at ($(heads.north)+(0,0.23)$) {selected $1{\times}1$ heads};

  \draw[queryarrow] (task.east) |- (posehead.west);
  \draw[queryarrow] (posehead.east) -- (poseout.west);
  \draw[taparrow] ($(task.south east)+(-0.23,0)$) |- (decoders.west);
  \node[label, text=gsblue!85!black, anchor=east] at (10.32,2.82)
    {task-layer patch taps};
  \draw[taparrow] (decoders.east) -- (heads.west);

  \node[label, text width=5.8cm] at (13.75,0.35)
    {Dense heads restore image space; global mixing stays in the task Transformer.};
\end{tikzpicture}
}
  \caption{\textbf{大域クエリと機械的な密予測デコーディング。}  \DINO{} patch tokenを
  射影し、2次元rotary encodingを付与する。1つの学習可能なcourt queryには位置回転を
  適用せず、task Transformer内ですべてのpatchに対して大域アテンションを行う。その最終状態は
  pose10Dを予測する。選択されたpatch tapは交換可能な軽量デコーダへ渡され、\KP、ライン、
  セグメンテーションの各直接予測ヘッドへ接続される。}
  \label{fig:model}
\end{figure}

\subsection{厳密なpatch-token境界}

拡張後の画像 \(I\in\R^{H\times W\times 3}\) が与えられると、その右端と下端を
replicate paddingして \DINO{} のpatch倍数に揃える。backboneはpatch tokenとspecial tokenを
含むsequenceを返す。adapterは \(H_pW_p\) 個のpatch tokenのみを受理し、token数、patch grid、
dtype、device、および有限性が一致することを検証する。CLS tokenとregister tokenは除外し、
task encoder内で曖昧性のない唯一の大域トークンが、自身の学習可能なcourt queryとなるようにする。

patch feature \(\mathbf{x}_{ij}\) はtask幅 \(d\) に射影される。各attention layerにおいて、
2次元rotary position encodingが行と列の関数としてpatchのquery/key pairを回転させる
\citep{su2021roformer}。先頭に付加した学習可能token \(\mathbf{q}_0\) は回転させない。
\begin{equation}
  \mathbf{Z}^{(0)}=
  [\mathbf{q}_0,\,\mathbf{P}\mathbf{x}_{00},\ldots,
   \mathbf{P}\mathbf{x}_{H_p-1,W_p-1}],
  \qquad
  \operatorname{RoPE}_{2D}(\mathbf{q}_0)=\mathbf{q}_0.
  \label{eq:query_input}
\end{equation}
すべてのquery tokenとpatch tokenが同じself-attention layerに参加する。したがって、最終的な
query state \(\mathbf{q}_L\) はコートの任意の可視部分から情報を集約できる一方で、patch stateは
自身の空間的identityを保持する。

\subsection{単一のカメラ仮説}
\label{sec:pose10d}

対象ビューはまず、カメラ側エンドに適用する行列式が $+1$ の回転行列によって正準化される。
\begin{equation}
  \T{\Canon}{\Cam}=\mathbf{S}\T{\Court_i}{\Cam},
  \qquad \mathbf{S}\in\{\mathbf{I},\mathbf{R}_z(\pi)\}.
  \label{eq:canonical_pose}
\end{equation}
物理的な \KP の順序とcanonicalizationの選択は、正とする情報として格納される。モデルアダプタは
permutationを推論せず、密予測教師と姿勢教師の不一致を棄却する。

MLPは \(\mathbf{q}_L\) を
\begin{equation}
  \widehat{\mathbf{p}}=
  [\widehat C_x,\widehat C_y,\widehat C_z,
   \widehat a_{11},\widehat a_{12},\widehat a_{13},
   \widehat a_{21},\widehat a_{22},\widehat a_{23},
   \widehat\ell_f]\in\R^{10}.
  \label{eq:pose10d}
\end{equation}
へ写す。最初の3値は、正準コートのメートル座標におけるカメラ中心
\(\widehat{\mathbf{C}}\) である。続く2つの未処理の行ベクトルは、次のように直交化する。
\begin{align}
  \mathbf{r}_1&=\frac{\mathbf{a}_1}{\|\mathbf{a}_1\|_2}, &
  \bar{\mathbf{r}}_2&=\mathbf{a}_2-
      (\mathbf{r}_1^{\top}\mathbf{a}_2)\mathbf{r}_1, \\
  \mathbf{r}_2&=\frac{\bar{\mathbf{r}}_2}{\|\bar{\mathbf{r}}_2\|_2}, &
  \mathbf{r}_3&=\mathbf{r}_1\times\mathbf{r}_2, \\
  \widehat{\mathbf{R}}_{\Canon\leftarrow\Cam}
    &= [\mathbf{r}_1^{\top};\mathbf{r}_2^{\top};\mathbf{r}_3^{\top}], &
  \widehat f&=\exp(\widehat\ell_f).
  \label{eq:rotation_decode}
\end{align}
この連続表現はquaternionの正規化と符号の曖昧性を回避し、いずれかの基底ベクトルが
退化した場合には棄却される\citep{zhou2019continuity}。焦点距離は有限かつ厳密に正である。
主点は予測せず、augmentation-awareなtarget intrinsicsから取得する。

\subsection{軽量な密予測ヘッド}

task Transformerから選択したpatch stateを、次の3つの制御されたデコーダ族のいずれかへ
渡す。
\begin{itemize}
  \item \emph{linear}: 各位置での射影と、それに続くresize。
  \item \emph{progressive}: depthwise-separable refinementを伴う段階的upsampling。
  \item \emph{DPT}: multi-tap reassemblyとfeature fusion。
\end{itemize}
デコーダは画像解像度のfeature mapを生成する。独立した \(1\times1\) ヘッドが、14個のsingleton
キーポイントlogit、1枚のline-mask logit map、およびcourt segmentation logitを出力する。
デコーダ幅、stage、fusion mode、およびtapは型付けされたablation変数だが、出力contractは
変化しない。この設計により、不透明で高容量なtaskヘッドが大域幾何を再構成するのではなく、
共有token encoder内に大域幾何を保持できるかを検証する。

直接目的関数は引き続き不可欠である。
\begin{align}
  \mathcal{L}_{\mathrm{dense}}
   &=w_{\mathrm{kp}}\mathcal{L}_{\mathrm{focalBCE}}
    +w_{\mathrm{line}}\mathcal{L}_{\mathrm{BCE+Dice}}
    +w_{\mathrm{seg}}\mathcal{L}_{\mathrm{CE+Dice}}, \\
  \mathcal{L}_{\mathrm{pose}}
   &=w_t\|\widehat{\mathbf{C}}-\mathbf{C}\|_2^2
    +w_R\,d_{\SOthree}(\widehat{\mathbf{R}},\mathbf{R})
    +w_f(\widehat\ell_f-\ell_f)^2,
  \label{eq:direct_losses}
\end{align}
ここで \(d_{\SOthree}\) は測地的回転距離である。次節の補助目的関数は、それまで独立していた
これらの分岐を接続するが、各分岐の正解損失を置き換えることはない。
